\section{Definición del Problema}

\noindent El problema estudiado corresponde al Multi-Source Capacitated Facility Location Problem with Customer Incompatibilities (MS-CFLP-CI), una variante del problema de localización de instalaciones capacitadas con incompatibilidades entre clientes. En términos generales, el problema consiste en decidir qué bodegas, instalaciones o centros de distribución deben abrirse y cómo debe distribuirse la demanda de los clientes entre ellas. Cada bodega posee una capacidad máxima y un costo fijo de apertura, mientras que atender a un cliente desde una bodega genera un costo operacional asociado. El objetivo es encontrar una solución de costo mínimo, considerando tanto los costos de apertura como los costos de operación necesarios para satisfacer completamente la demanda de todos los clientes \cite{klose2005facility,maia2023metaheuristic,ceschia2024multineighborhood}.\\

\noindent La característica principal del MS-CFLP-CI es que un cliente puede recibir su demanda desde más de una bodega. Es decir, la demanda de un cliente puede dividirse entre distintas instalaciones, siempre que la suma total recibida satisfaga completamente su requerimiento. Esta condición diferencia a la variante multi-source de la variante single-source, donde cada cliente debe ser atendido por una única instalación. La posibilidad de dividir la demanda entrega mayor flexibilidad al modelo, ya que permite aprovechar mejor las capacidades disponibles y combinar alternativas de menor costo \cite{mansini2022kernel,ceschia2024multineighborhood}.\\

\noindent Además de las restricciones clásicas de capacidad y demanda, el MS-CFLP-CI incorpora incompatibilidades entre clientes. Estas incompatibilidades se representan mediante pares de clientes que no pueden ser atendidos por una misma bodega. En la práctica, esta situación puede aparecer cuando ciertos productos, servicios o requerimientos no deben compartir una misma instalación por razones sanitarias, técnicas, operativas o de seguridad. Por lo tanto, una solución factible debe asegurar que cada cliente reciba toda su demanda, que ninguna bodega exceda su capacidad y que ningún par de clientes incompatibles sea servido desde la misma instalación \cite{maia2023metaheuristic,mansini2022kernel,ceschia2024multineighborhood}.\\

\noindent Las decisiones principales del problema se pueden describir en tres grupos. Primero, se debe determinar qué bodegas serán abiertas, lo que implica asumir sus respectivos costos fijos. Segundo, se debe decidir cuánta demanda de cada cliente será enviada desde cada bodega abierta. Tercero, se debe verificar que las asignaciones realizadas no generen conflictos de incompatibilidad. En la variante MS-CFLP-CI, esta última verificación es especialmente importante, ya que un cliente queda asociado a cada bodega desde la cual recibe una cantidad positiva de productos \cite{mansini2022kernel,ceschia2024multineighborhood}.\\

\noindent Desde el punto de vista computacional, el MS-CFLP-CI es un problema de alta dificultad, ya que hereda la complejidad del CFLP clásico, reconocido como un problema NP-Hard, e incorpora restricciones adicionales de incompatibilidad entre clientes. Además, el tamaño del espacio de búsqueda crece rápidamente: para $m$ bodegas existen $2^m$ posibles combinaciones de apertura, y para $n$ clientes la asignación multi-source puede considerar múltiples subconjuntos de bodegas para cada cliente. Esto genera un espacio combinatorio muy grande, especialmente en instancias con muchas instalaciones, clientes e incompatibilidades \cite{klose2005facility,maia2023metaheuristic,ceschia2024multineighborhood}.\\

\noindent El MS-CFLP-CI también puede relacionarse con otras variantes de problemas de localización, como el Uncapacitated Facility Location Problem (UFLP), donde no se consideran límites de capacidad; el $p$-Median Problem, que busca ubicar un número fijo de instalaciones minimizando distancias o costos de asignación; y el Robust Facility Location Problem, que incorpora incertidumbre en parámetros como demanda, costos o disponibilidad. A diferencia de estas variantes, el MS-CFLP-CI combina capacidades limitadas, asignación divisible de demanda e incompatibilidades entre clientes, por lo que requiere enfoques de solución más especializados, incluyendo modelos matemáticos, heurísticas, metaheurísticas y matheurísticas \cite{fischetti2016benders,maia2023metaheuristic,mansini2022kernel,ceschia2024multineighborhood}.